\subsection{The Nullstellensatz}

PROPOSITION 1. Let A be a finitely generated algebra over a field \(k\) and \(L\) the algebraic closure of \(k\) .

\begin{enumerate}
\def\labelenumi{(\roman{enumi})}
\item
  If \(\mathbf{A} \neq \{0\}\) , there exists a \(k\) -homomorphism from \(\mathbf{A}\) to \(\mathbf{L}\) .
\item
  Let \(f_{1}, f_{2}\) be two \(k\) -homomorphisms from \(\mathbf{A}\) to \(\mathbf{L}\) . For \(f_{1}\) and \(f_{2}\) to have the same kernel, it is necessary and sufficient that there exist a \(k\) -automorphism \(s\) of \(\mathbf{L}\) such that \(f_{2} = s \circ f_{1}\) .
\item
  Let \(\mathfrak{a}\) be an ideal of \(\mathbf{A}\) . For \(\mathfrak{a}\) to be maximal, it is necessary and sufficient that it be the kernel of a \(k\) -homomorphism from \(\mathbf{A}\) to \(\mathbf{L}\) .
\item
  For an element \(x\) of \(\mathbf{A}\) to be such that \(f(x) = 0\) for every \(k\) -homomorphism \(f\) from \(\mathbf{A}\) to \(\mathbf{L}\) , it is necessary and sufficient that \(x\) be nilpotent.
\end{enumerate}

Assertion (i) follows from no. 1, Corollary 3 to Theorem 1 applied replacing A by k, B by A, b by the unit element of B and f by the canonical injection of k into L.

If \(f\) is a \(k\) -homomorphism from \(\mathbf{A}\) to \(\mathbf{L}\) , \(f(\mathbf{A})\) is a subring of \(\mathbf{L}\) containing \(k\) ; as \(\mathbf{L}\) is an algebraic extension of \(k\) , \(f(\mathbf{A})\) is a field (Algebra, Chapter V, § 3, no. 2, Proposition 3) and, if \(\mathfrak{a}\) is the kernel of \(f\) , \(\mathbf{A} / \mathfrak{a}\) , isomorphic to \(f(\mathbf{A})\) , is therefore a field, which proves that \(\mathfrak{a}\) is maximal. Conversely, if \(\mathfrak{a}\) is a maximal ideal of \(\mathbf{A}\) , it follows from (i) that there exists a \(k\) -homomorphism from \(\mathbf{A} / \mathfrak{a}\) to \(\mathbf{L}\) and hence a \(k\) -homomorphism of \(\mathbf{A}\) to \(\mathbf{L}\) whose kernel \(\mathfrak{b}\) contains \(\mathfrak{a}\) ; but as \(\mathfrak{a}\) is maximal, \(\mathfrak{b} = \mathfrak{a}\) ; this proves (iii).

We now prove (ii). If \(s\) is a \(k\) -automorphism of \(L\) such that \(f_2 = s \circ f_1\) , clearly \(f_1\) and \(f_2\) have the same kernel. Conversely, suppose that \(f_1\) and \(f_2\) have the same kernel; then there exists a \(k\) -isomorphism \(s_0\) of the field \(f_1(A)\) onto the field \(f_2(A)\) such that \(f_2 = s_0 \circ f_1\) ; but by Algebra, Chapter V, § 6, no. 3, Proposition 7, \(s_0\) extends to a \(k\) -automorphism \(s\) of \(L\) and hence \(f_2 = s \circ f_1\) .

Finally, if \(x \in \mathbf{A}\) is such that \(x^n = 0\) , for every \(k\) -homomorphism \(f\) from \(\mathbf{A}\) to \(\mathbf{L}\) , \((f(x))^n = f(x^n) = 0\) and hence \(f(x) = 0\) since \(\mathbf{L}\) is a field. Conversely, suppose that \(x \in \mathbf{A}\) is not nilpotent; then \(\mathbf{A}[x^{-1}]\) is a finitely generated A-algebra (and therefore a finitely generated \(k\) -algebra) not reduced to 0 (Chapter II, §2, no. 1, Remark 3) and hence there exists a \(k\) -homomorphism \(g\) from \(\mathbf{A}[x^{-1}]\) to \(\mathbf{L}\) by (i). If \(j: \mathbf{A} \to \mathbf{A}[x^{-1}]\) is the canonical homomorphism, \(f = g \circ j\) is a \(k\) -homomorphism from \(\mathbf{A}\) to \(\mathbf{L}\) and \(f(x)g(1/x) = g(x/1)g(1/x) = g(1) = 1\) , whence \(f(x) \neq 0\) .

Let \(k\) be a field and \(\mathbf{L}\) an extension field of \(k\) ; an element \(\mathbf{x} = (x_1, \ldots, x_n)\) of \(\mathbf{L}^n\) is called a zero in \(\mathbf{L}^n\) of an ideal \(\mathfrak{r}\) of the polynomial ring \(k[\mathbf{X}_1, \ldots, \mathbf{X}_n]\) if

\[
\mathrm{P} (\mathbf {x}) = \mathrm{P} (x _ {1}, \dots , x _ {n}) = 0
\]

for all \(\mathbf{P} \in \mathfrak{r}\) .

LEMMA 1. Let A be a finitely generated algebra over a field \(k\) , \((a_i)_{1 \leqslant i \leqslant n}\) a system of generators of this algebra and \(\mathfrak{r}\) the ideal of algebraic relations between the \(a_i\) with coefficients in \(k\) (Algebra, Chapter IV, §2, no. 1). For every extension field \(\mathbf{L}\) of \(k\) , the mapping \(f \mapsto (f(a_i))_{1 \leqslant i \leqslant n}\) is a bijection of the set of \(k\) -homomorphisms from A to L onto the set of zeros of \(\mathfrak{r}\) in \(L^n\) .

There exists a unique \(k\) -algebra homomorphism \(h\) of \(k[\mathbf{X}_1, \ldots, \mathbf{X}_n]\) onto \(\mathbf{A}\) such that \(h(\mathbf{X}_i) = a_i\) for \(1 \leqslant i \leqslant n\) and by definition \(\mathfrak{r}\) is the kernel of \(h\) . The mapping \(f \mapsto f \circ h\) is a bijection of the set of \(k\) -homomorphisms from \(A\) to \(L\) onto the set of \(k\) -homomorphisms from \(k[X_1, \ldots, X_n]\) to \(L\) which are zero on \(r\) . For every polynomial \(P \in k[X_1, \ldots, X_n]\) and every element \(\mathbf{x} = (x_1, \ldots, x_n) \in L^n\) we write \(h_{\mathbf{x}}(P) = P(\mathbf{x})\) ; then the mapping \(\mathbf{x} \mapsto h_{\mathbf{x}}\) is a bijection of \(L^n\) onto the set of \(k\) -homomorphisms from \(k[X_1, \ldots, X_n]\) to \(L\) (such a homomorphism being determined by its values at the \(X_i\) ( \(1 \leqslant i \leqslant n\) )); to say that \(h_{\mathbf{x}}\) is zero on \(r\) means that \(\mathbf{x}\) is a zero of \(r\) in \(L^n\) , whence the lemma.

If Proposition 1 is applied to the algebra \(\mathbf{A} = k[\mathbf{X}_1, \ldots, \mathbf{X}_n] / \mathfrak{r}\) , where \(\mathfrak{r}\) is an ideal of \(k[\mathbf{X}_1, \ldots, \mathbf{X}_n]\) distinct from the whole ring, we obtain by Lemma 1 the following statement:

PROPOSITION 2 (Hilbert's Nullstellensatz). Let \(k\) be a field and \(L\) an algebraic closure of \(k\) .

\begin{enumerate}
\def\labelenumi{(\roman{enumi})}
\item
  Every ideal \(\mathfrak{r}\) of \(k[\mathbf{X}_1, \ldots, \mathbf{X}_n]\) not containing 1 admits at least one zero in \(\mathbf{L}^n\) .
\item
  Let \(\mathbf{x} = (x_1, \ldots, x_n)\) , \(\mathbf{y} = (y_1, \ldots, y_n)\) be two elements of \(\mathbf{L}^n\) ; for the set of polynomials of \(k[\mathbf{X}_1, \ldots, \mathbf{X}_n]\) zero at \(\mathbf{x}\) to be identical with the set of polynomials of \(k[\mathbf{X}_1, \ldots, \mathbf{X}_n]\) zero at \(\mathbf{y}\) , it is necessary and sufficient that there exists a \(k\) -automorphism \(s\) of \(\mathbf{L}\) such that \(y_i = s(x_i)\) for \(1 \leqslant i \leqslant n\) .
\item
  For an ideal \(\mathfrak{a}\) of \(k[\mathbf{X}_1, \ldots, \mathbf{X}_n]\) to be maximal, it is necessary and sufficient that there exist an \(\mathbf{x}\) in \(\mathbf{L}^n\) such that \(\mathfrak{a}\) is the set of polynomials of \(k[\mathbf{X}_1, \ldots, \mathbf{X}_n]\) zero at \(\mathbf{x}\) .
\item
  For a polynomial \(\mathbf{Q}\) of \(k[\mathbf{X}_1, \ldots, \mathbf{X}_n]\) to be zero on the set of zeros in \(\mathbf{L}^n of\) an ideal \(\mathfrak{r}\) of \(k[\mathbf{X}_1, \ldots, \mathbf{X}_n]\) , it is necessary and sufficient that there exist an integer \(m > 0\) such that \(\mathbf{Q}^m \in \mathfrak{r}\) .
\end{enumerate}

